\color{unidarkgreen}{\subsection{Логика сигнализации}}
\color{black}

\begin{enumerate}[label=\arabic{section}.\arabic{subsection}.\arabic*,  align=left, labelsep=4pt, leftmargin=0pt, itemindent=57pt]

\item
Устройство формирует выходное воздействие на сигнализацию \textit{«Сигн. Сраб.(фикс)»} и \textit{«Откл. защ.
(фикс)»} по схеме «ИЛИ» в следующих случаях:
\begin{itemize}
\item Отключении ШР;
\item Срабатывании ГЗ ШР;
\item Срабатывании ГЗ ШР сигн;
\item Срабатывании отсечного клапана; 
\item Срабатывании ЗПО;
\item Срабатывании ЗНР;
\item Срабатывании сигн. КИВ;
\item Срабатывании КОН.
\end{itemize}

\item
Устройство формирует выходное воздействие на сигнализацию \textit{«Сигн. неиспр.»} и \textit{«Сигн. неиспр.
(фикс)»} по схеме «ИЛИ» в следующих случаях:
\begin{itemize}
\item Неисправности цепей тока ДТЗ;
\item Неисправности цепей ГЗ ШР;
\item Неисправность ТЗ ШР;
\item Максимальный уровень масла;
\item Минимальный уровень масла;
\item Отказ сист. охл.;
\item Неисправность охлаждения;
\item Неисправность ОТ терминала;
\item Неисправность КИВ.
\end{itemize}

\item
Сброс фиксации производится с помощью \textit{«Сброс»}.

\begin{figure}[H]
\centering
\includegraphics[width=0.7\textwidth,height=0.7\textheight,keepaspectratio]{\fbpath/91 Логика сигнализации/img/ls.pdf}
\captionof{figure}{Функциональный блок логики сигнализации}\label{pict_ls:LS}
\end{figure}

\end{enumerate}